\documentclass[10pt,a4paper]{amsart}
\usepackage[margin=19mm]{geometry}
\usepackage{amsmath,amssymb,mathtools}
\usepackage[scr=rsfs,cal=euler]{mathalfa}
\usepackage{xfrac}
\usepackage[unicode,hidelinks,bookmarks=false]{hyperref}
\newcommand{\ie}{i.e.\ }
\newcommand{\cf}{cf.\ }
\newcommand{\resp}{respectively\ }
\newcommand{\Subsection}{\subsection}
\providecommand{\nobreakdash}{\nobreak\hbox{-}}
\setlength{\emergencystretch}{2em}
\multlinegap0pt
\usepackage{amssymb}
\usepackage{float}
\usepackage{tikz}
\newtheorem{theorem}{Theorem}
\newcommand{\wt}{{\texttt{wt}}}
\title{Holonomic Sequences Arising from Laplace Transforms of Generalized Second Order Sequences}
\author{Rigoberto Fl\'orez, Robinson Higuita-Diaz, Jos\'e L. Ram\'{\i}rez}
\date{Source: arXiv:2607.22955v1 (CC BY 4.0)}
\begin{document}
\maketitle

\noindent\emph{Source and licence.} Holonomic Sequences Arising from Laplace Transforms of Generalized Second Order Sequences. Version arXiv:2607.22955v1 (CC BY 4.0), distributed by arXiv under the Creative Commons Attribution 4.0 International licence, http://creativecommons.org/licenses/by/4.0/, as recorded in the arXiv licence metadata for this exact version and shown on the abstract page https://arxiv.org/abs/2607.22955v1. Full licence notice: \texttt{licenses/arxiv-2607.22955-CC-BY-4.0.txt}.\par\smallskip

\noindent\textit{Editorial scope.} Counted unit: the article's theorem thm:wn-recurrence (source0.tex lines 775--787). The article's complete original proof of it is delivered with it (source0.tex lines 788--856, one \texttt{proof} environment of 69 lines). No mathematics is written, repaired or re-derived: the statement and the proof are the article's own lines, in the article's own notation, and no retained line is substituted or re-flowed.  No further source-local context is needed: every cross-reference of every retained line resolves inside the two delivered spans.  Nothing was omitted from any retained span, so the package carries no omission record.  No retained line contains a citation, so the wrapper carries no bibliography and no bibliography entry.  No retained line contains a figure environment, an image-inclusion command or drawing code, so no image is delivered and the package declares no asset dependency.  Editorial formatting only, in addition: an AMS \texttt{amsart} preamble that restates the article's own declarations and macros used by the retained lines, namely the environments \texttt{theorem} on separate counters, together with the standard packages those lines need.  The retained environments print in this document as Theorem 1 in order of appearance, whereas the article prints the same items with the numbering of its own full document.  The complete native source of the article as distributed in the arXiv e-print for this exact version is bundled unchanged as \texttt{sources/205987/source0.tex}.
\begin{theorem}\label{thm:wn-recurrence}
The sequence $(w_n)_{n\ge0}$ satisfies, for $n\ge 4$,
\begin{equation}\label{eq:wn-rec}
w_n=(b+an)\,w_{n-1}-ab(n-1)\,w_{n-2}+\bigl(bc-ac(n-2)\bigr)\,w_{n-3}+c^2\,w_{n-4},
\end{equation}
with initial values
\[
w_0=1,\quad 
w_1=a+b,\quad
w_2=2a^2+2ab+b^2+c,\quad
w_3=6a^3+6a^2b+3ab^2+b^3+2ac+2bc.
\]
\end{theorem}

\begin{proof}
Let $\mathcal{Q}_n$ be the set of tilings in $\mathcal{C}_n$ with no type-$A$ squares, and let $q_n$ be their total weight. Such
tilings contain only type-$B$ squares and dominoes. Considering
the last tile gives
\begin{equation}\label{eq:gn-recurrence}
q_n=b\,q_{n-1}+c\,q_{n-2},
\qquad n\geq 2.
\end{equation}
Therefore $w_n-q_n$ is the total weight of the tilings in $\mathcal{C}_n$
containing at least one type-$A$ square. 

We first prove
\begin{equation}\label{eq:wn-positive-recurrence}
an\,w_{n-1}
=
w_n-q_n+c\bigl(w_{n-2}-q_{n-2}\bigr),
\qquad n\geq 2.
\end{equation}
Consider pairs $(P,\gamma)$, where $P\in\mathcal{C}_{n-1}$ and
$\gamma$ is one of the $n$ vertical grid lines of the board, including the two boundary lines. We assign weight $a\wt(P)$ to each such pair. Since there are $n$ choices for $\gamma$, the total weight of all
these pairs is $an\,w_{n-1}$.


Suppose that $P$ contains exactly $k$ type-$A$ squares. If $\gamma$ lies between two tiles or is a boundary line, insert a type-$A$ square
at $\gamma$ and give it color $k+1$.   This gives a tiling in $\mathcal{C}_n$ containing at least one type-$A$ square, and the
construction is weight preserving.

If $\gamma$ passes through a domino, replace that domino by a type-$A$
square of color $k+1$. This gives a tiling $F\in\mathcal{C}_{n-2}$
containing at least one type-$A$ square. Moreover, $a\,\wt(P)=c\,\wt(F)$, since a domino of weight $c$ has been replaced by a type-$A$ square
of weight $a$. 

Both constructions are reversible. In the first case, remove the
type-$A$ square with largest color and mark the grid line created by
its removal. In the second case, replace the type-$A$ square with
largest color by a domino and mark its central grid line.


Thus the first case contributes $w_n-q_n$, while the second contributes
$c\bigl(w_{n-2}-q_{n-2}\bigr)$. Since the two cases partition all
pairs $(P,\gamma)$, \eqref{eq:wn-positive-recurrence} follows.

Now define
\[
r_n:=w_n+c\,w_{n-2}-an\,w_{n-1}.
\]
By~\eqref{eq:wn-positive-recurrence}, $r_n=q_n+c\,q_{n-2}$. Using \eqref{eq:gn-recurrence}, for $n\geq 4$, we have
\begin{align*}
r_n
&=q_n+c\,q_{n-2}\\
&=b\,q_{n-1}+c\,q_{n-2}
  +bc\,q_{n-3}+c^2q_{n-4}\\
&=b\bigl(q_{n-1}+c\,q_{n-3}\bigr)
  +c\bigl(q_{n-2}+c\,q_{n-4}\bigr)\\
&=b\,r_{n-1}+c\,r_{n-2}.
\end{align*}
Substituting the definition of $r_n$ gives
\begin{multline*}
w_n+c\,w_{n-2}-an\,w_{n-1}
\\=
b\Bigl(
w_{n-1}+c\,w_{n-3}-a(n-1)w_{n-2}
\Bigr)+
c\Bigl(
w_{n-2}+c\,w_{n-4}-a(n-2)w_{n-3}
\Bigr).
\end{multline*}
Canceling $c\,w_{n-2}$ and collecting terms proves \eqref{eq:wn-rec}. The initial values follow by direct inspection of the tilings of lengths $0,1,2$, and $3$.
\end{proof}

\end{document}
